\documentclass{article}
\usepackage{amsmath,amssymb,mathtools}
\usepackage[margin=0.7in]{geometry}
\allowdisplaybreaks
\begin{document}

\section*{DRO location--inventory model: formulation (10) with the full epigraph (9)}
\begin{equation}
\begin{aligned}
\min_{\substack{\boldsymbol{o},\boldsymbol{z},\boldsymbol{\alpha},\boldsymbol{\phi}\\
\boldsymbol{\gamma},\boldsymbol{\lambda},\boldsymbol{\tau},\delta}}
\Biggl\{&\sum_{i\in I}f_i o_i+
\sum_{i\in I}\sum_{t\in\mathcal T}c_t^a z_{t,i}
+\sum_{i\in I}\sum_{t\in\mathcal T}\sum_{l\in L}
 \left(\mu^d_{t,i,l}\alpha_{t,i,l}+\mu^M_{t,i,l}\phi_{t,i,l}\right)\\
&+\sum_{i\in I}\sum_{l\in L}
 \left(\mu^\rho_{i,l}\gamma_{i,l}+\mu^q_{i,l}\lambda_{i,l}\right)
+\sum_{(i,j)\in\mathcal A}\sum_{l\in L}\mu^V_{i,j,l}\tau_{i,j,l}+\delta\Biggr\}
\end{aligned}\tag{10a}
\end{equation}
\begin{align}
\text{s.t.}\quad &(1\mathrm b)\text{--}(1\mathrm c),\tag{10b}\\
&\delta\geq H(\boldsymbol z,\boldsymbol\alpha,\boldsymbol\phi,
\boldsymbol\gamma,\boldsymbol\lambda,\boldsymbol\tau).\tag{10c}
\end{align}
The inherited first-stage rows are
\begin{align}
&\sum_{t\in\mathcal T}s_t z_{t,i}\leq S_i o_i,
&&\forall i\in I,\tag{1b}\\
&o_i\in\{0,1\},\quad z_{i,t\geq0},
&&\forall i\in I,\ t\in\mathcal T.\tag{1c}
\end{align}
\noindent\emph{Printed notation in (1c).} The page image prints the malformed subscript
$z_{i,t\geq0}$, retained here. The defining paragraph on page 10 separately specifies
$z_{t,i}$ as a nonnegative continuous variable:
\[
z_{t,i}\geq0\qquad(t\in\mathcal T,\ i\in\mathcal N).
\]
No identification of an additional variable $z_{i,t}$ is introduced.

\subsection*{Sets, fixed data, and variable roles}
The directed network is $G(\mathcal N,\mathcal A)$. The set $\mathcal N$ consists of nodes,
$\mathcal A$ of directed road arcs, $I\subseteq\mathcal N$ of potential warehouse sites,
$\mathcal T$ of relief-item types, and $L$ of disaster levels. The problem-description
paragraph assumes that candidate locations and demand nodes are the same; Table 1
nevertheless states $I\subseteq\mathcal N$, and the displayed model uses $I$ as retained above.

The fixed nonnegative parameters $f_i,S_i,s_t,c_t^a,c_t^p,c_t^u,c_t^h,c^t_{i,j}$ are,
respectively, facility-opening cost, warehouse storage capacity, unit storage requirement,
unit acquisition cost, unit post-disaster procurement cost, unit unmet-demand penalty,
unit holding cost, and unit transportation cost. Their indices are $i\in\mathcal N$,
$t\in\mathcal T$, and $(i,j)\in\mathcal A$, as applicable. The parameter $v_t$ is arc
capacity required per unit of item $t$; $\hat V_{i,j}$ is nominal arc capacity.
The source assumes no limit on supplies available for first-stage prepositioning.
It describes post-disaster procurement as limited and more expensive, without specifying
a general numerical cost ratio.

The first-stage decision $o_i$ indicates whether a facility is open, and $z_{t,i}$ is its
prepositioned quantity of item $t$. The five families
$\alpha_{t,i,l},\phi_{t,i,l},\gamma_{i,l},\lambda_{i,l},\tau_{i,j,l}$ are moment-dual
variables, not physical inventory or flow variables. Appendix B specifies them as
unrestricted in sign, for $t\in\mathcal T$, $i\in I$, $l\in L$, and $(i,j)\in\mathcal A$
as appropriate. The scalar $\delta$ is the full recourse epigraph variable; no
nonnegativity restriction on $\delta$ is printed.

\subsection*{Uncertainty, support, and ambiguity}
The random vector is
\[
\xi=[\boldsymbol q,\boldsymbol d,\boldsymbol\rho,\boldsymbol M,\boldsymbol V]^\top.
\]
For $i\in I$, $l\in L$, $q_{i,l}$ is binary and equals one if a level-$l$ disaster occurs
at node $i$, and zero otherwise. The random quantity $d_{t,i,l}$ is demand for item $t$
following that disaster; $\rho_{i,l}$ is the fraction of prepositioned supplies remaining
usable; $M_{t,i,l}$ is the maximum post-disaster procurement quantity; and $V_{i,j,l}$ is
arc capacity following a level-$l$ disaster at either endpoint. These are uncertain
inputs to recourse, but their realizations are chosen in the worst-case maximization
underlying $H$. Means and support endpoints are fixed data, assumed known or approximable.

The source defines
\[
\mathcal R=\mathcal R_q\times\mathcal R_d\times\mathcal R_\rho
\times\mathcal R_M\times\mathcal R_V,
\]
\begin{align}
\mathcal R_q&:=\{0,1\}^{N},&
\mathcal R_d&:=\left\{\boldsymbol d\geq0:
\underline d_{t,i,l}\leq d_{t,i,l}\leq\overline d_{t,i.l},\ \forall(t,i,l)\right\},\tag{3a}\\
\mathcal R_\rho&:=\left\{\boldsymbol\rho\geq0:
\underline\rho_{i,l}\leq\rho_{i,l}\leq\overline\rho_{i,l},\ \forall(i,l)\right\},\tag{3b}\\
\mathcal R_M&:=\left\{\boldsymbol M\geq0:
\underline M_{t,i,l}\leq M_{t,i,l}\leq\overline M_{t,i,l},\ \forall(t,i)\right\},\tag{3c}\\
\mathcal R_V&:=\left\{\boldsymbol V\geq0:
\underline V_{i,j,l}\leq V_{i,j,l}\leq\overline V_{i,j,l},\
\forall(i,j)\in\mathcal A,\ \forall l\right\}.\tag{3d}
\end{align}
Here lower and upper decorations denote support endpoints, not random realizations.
The supplied text of (3a) uses exponent $N$ and the punctuation $i.l$ in the demand
upper endpoint; (3c) quantifies $(t,i)$ although $l$ occurs in its expression. These
printed scopes are retained. The component indexing in the definition paragraphs is
$t\in\mathcal T$, $i\in\mathcal N$ for $d,M$, $i\in I$ for $q,\rho$, $l\in L$, and
$(i,j)\in\mathcal A$ for $V$. No replacement dimension for $\mathcal R_q$ is supplied here.

Let $\mathcal P=\mathcal P(\mathcal R)$ be the probability measures on
$(\mathcal R,\mathcal B)$, with $\mathcal B$ the Borel $\sigma$-field, and let
$\mathbb E_{\mathbb P}$ denote expectation under $\mathbb P$. Then
\[
\mu:=\mathbb E_{\mathbb P}[\xi]
=[\boldsymbol\mu^q,\boldsymbol\mu^d,\boldsymbol\mu^\rho,
\boldsymbol\mu^M,\boldsymbol\mu^V]^\top,
\]
\begin{equation}
\mathcal F(\mathcal R,\mu):=
\left\{\mathbb P\in\mathcal P(\mathcal R)\mid
\mathbb E_{\mathbb P}[\xi]=\mu\right\}.\tag{4}
\end{equation}
The component moment and normalization equalities used in Appendix B are
\begin{align}
\int_{\mathcal R}d_{t,i,l}\,d\mathbb P&=\mu^d_{t,i,l},
&&\forall i\in I,\ t\in\mathcal T,\ l\in L,\tag{B.1b}\\
\int_{\mathcal R}M_{t,i,l}\,d\mathbb P&=\mu^M_{t,i,l},
&&\forall i\in I,\ t\in\mathcal T,\ l\in L,\tag{B.1c}\\
\int_{\mathcal R}\rho_{i,l}\,d\mathbb P&=\mu^\rho_{i,l},
&&\forall i\in I,\ l\in L,\tag{B.1d}\\
\int_{\mathcal R}V_{i,j,l}\,d\mathbb P&=\mu^V_{i,j,l},
&&\forall(i,j\in\mathcal A,\ l\in L,\tag{B.1e}\\
\int_{\mathcal R}q_{i,l}\,d\mathbb P&=\mu^q_{i,l},
&&\forall i\in I,\ l\in L,\tag{B.1f}\\
\int_{\mathcal R}d\mathbb P&=1.\tag{B.1g}
\end{align}
The unmatched parenthesis in the quantifier of (B.1e) is retained from the supplied
source text. The multipliers of these six families are, respectively,
$\alpha,\phi,\gamma,\tau,\lambda,\theta$, all unrestricted in sign. The normalization
multiplier $\theta$ belongs to the derivation, not to the decision list of (10).

\paragraph{Disaster exclusivity and optional counts.}
Page 10 restricts the computational study, rather than the general formulation, to
\[
\sum_{l\in L}q_{i,l}\leq1,\qquad\forall i\in I.
\]
Likewise, page 16 says that to prescribe the number of disasters of level $l$, one
\emph{can add}
\[
\sum_{i\in I}q_{i,l}=N_l.
\]
Neither restriction is included in the general (9b) selected here. The case-study
values and hurricane-count equalities are not imposed on this formulation.

\subsection*{Full worst-case recourse epigraph}
For fixed $(\boldsymbol z,\boldsymbol\alpha,\boldsymbol\phi,
\boldsymbol\gamma,\boldsymbol\lambda,\boldsymbol\tau)$, Proposition 2 defines
\begin{equation}
\begin{aligned}
&H(\boldsymbol z,\boldsymbol\alpha,\boldsymbol\phi,
\boldsymbol\gamma,\boldsymbol\lambda,\boldsymbol\tau)\\
&=\max_{\substack{
\boldsymbol\beta,\boldsymbol\Gamma,\boldsymbol\psi,\boldsymbol\varphi,
\boldsymbol a,\boldsymbol\Theta,\boldsymbol\kappa,\boldsymbol\varrho\\
\boldsymbol a',\boldsymbol\Theta',\boldsymbol\kappa',\boldsymbol\varrho',
\boldsymbol h,\boldsymbol g,\boldsymbol F,\boldsymbol\pi\\
\boldsymbol\eta,\boldsymbol\Phi,\boldsymbol\varpi,\boldsymbol b,\boldsymbol\Lambda}}
\Biggl\{\sum_{i\in I}\sum_{t\in\mathcal T}\sum_{l\in L}
\Bigl[\underline d_{t,i,l}(k_{t,i,l}-\alpha_{t,i,l})
+\Delta d_{t,i,l}(h_{t,i,l}-\alpha_{t,i,l}a_{t,i,l})\Bigr]\\
&\qquad+\sum_{i\in I}\sum_{t\in\mathcal T}z_{t,i}
\left(\sum_{l\in L}k_{t,i,l}-\beta_{t,i}\right)\\
&\qquad-\sum_{i\in I}\sum_{l\in L}
\Bigl[\underline\rho_{i,l}\left(\sum_{t\in\mathcal T}z_{t,i}k_{t,i,l}+\gamma_{i,l}\right)
+\Delta\rho_{i,l}\left(\sum_{t\in\mathcal T}z_{t,i}g_{t,i,l}+\Theta_{i,l}\gamma_{i,l}\right)\Bigr]\\
&\qquad+\sum_{i\in I}\sum_{t\in\mathcal T}\sum_{l\in L}
\Bigl[\underline M_{t,i,l}(F_{t,i,l}-\phi_{t,i,l})
+\Delta M_{t,i,l}(\pi_{t,i,l}-\phi_{t,i,l}\kappa_{t,i,l})\Bigr]\\
&\qquad+\sum_{(i,j)\in\mathcal A}\sum_{l\in L}
\Bigl[\underline V_{i,j,l}(\eta_{i,j,l}+\varpi_{i,j,l}-\tau_{i,j,l})
+\Delta V_{i,j,l}(\Phi_{i,j,l}+\Lambda_{i,j,l}-\varrho_{i,j,l}\tau_{i,j,l})\\
&\qquad+\sum_{(i,j)\in\mathcal A}
\Bigl[\left(\psi_{i,j}-\sum_{l\in L}\eta_{i,j,l}\right)\hat V_{i,j}
+\left(\varphi_{i,j}-\sum_{l\in L}\varpi_{i,j,l}\right)\hat V_{i,j}\Bigr]
-\sum_{i\in I}\sum_{l\in L}q_{i,l}\lambda_{i,l}\Biggr\}
\end{aligned}\tag{9a}
\end{equation}
\begin{equation}
\text{s.t.}\quad
\{(7\mathrm a)\text{--}(7\mathrm d)\},\quad
\boldsymbol q\in\{0,1\},\quad
(\boldsymbol a,\boldsymbol\Theta,\boldsymbol\kappa,\boldsymbol\varrho)\in\{0,1\},
\quad (\mathrm C.3a)\text{--}(\mathrm C.3n).\tag{9b}
\end{equation}
\noindent\emph{Printed omissions in (9a).} The displayed maximization subscript omits
$\boldsymbol q$ and $\boldsymbol k$, although they occur in its objective and defining
constraints. The opening square bracket following
$\sum_{(i,j)\in\mathcal A}\sum_{l\in L}$ has no matching closing square bracket in the
original page image. Its printed placement is retained; no additional nesting or
closing location is imposed. The nominal-capacity and $q\lambda$ sums are reproduced
on their separate printed line.

The binary families have component indices
\[
q_{i,l},\Theta_{i,l}\ (i\in I,l\in L),\qquad
 a_{t,i,l},\kappa_{t,i,l}\ (t\in\mathcal T,i\in I,l\in L),\qquad
\varrho_{i,j,l}\ ((i,j)\in\mathcal A,l\in L).
\]
The remaining inner variables are continuous auxiliary or recourse-dual variables;
no additional integrality is assigned to the primed product variables.
The dual variables $\beta_{t,i},\Gamma_{t,i},\psi_{i,j},\varphi_{i,j}$ correspond to
(2b), (2c), (2d), and (2e), respectively. In particular, the outer multiplier
$\phi_{t,i,l}$ is distinct from the arc-capacity dual $\varphi_{i,j}$.

\subsection*{Dual feasible region inherited by (9)}
\begin{align}
-\beta_{t,i}+\beta_{t,j}+v_t\psi_{i,j}+v_t\varphi_{i,j}
&\leq c^t_{i,j},
&&\forall(i,j)\in\mathcal A,\ \forall t\in\mathcal T,\tag{7a}\\
\beta_{t,i}+\Gamma_{t,i}&\leq c_t^p,
&&\forall i\in I,\ t\in\mathcal T,\tag{7b}\\
-c_t^h\leq\beta_{t,i}&\leq c_t^u,
&&\forall i\in I,\ t\in\mathcal T,\tag{7c}\\
\Gamma_{t,i}\leq0,\quad\psi_{i,j}&\leq0,\quad\varphi_{i,j}\leq0,
&&\forall(i,j)\in\mathcal A,\ t\in\mathcal T.\tag{7d}
\end{align}
The quantifier in (7d) is retained as printed. Page 16 calls this region bounded and
$\mathcal R$ finite. Those claims are not extra constraints: the displayed region has
unbounded negative directions in its nonpositive dual variables, and the interval
supports in (3) are not generally finite sets.

\subsection*{Endpoint construction and auxiliary meanings}
Appendix C replaces each interval realization by a lower endpoint plus its range
multiplied by a binary endpoint selector. Its stated premise is that, at an optimum
of the inner maximization with the outer variables fixed, each such realization can
be taken at an endpoint. The defining relations are
\begin{align*}
\Delta d_{t,i,l}&=\overline d_{t,i,l}-\underline d_{t,i,l},&
d_{t,i,l}&=\underline d_{t,i,l}+a_{t,i,l}\Delta d_{t,i,l},\\
\Delta\rho_{i,l}&=\overline\rho_{i,l}-\underline\rho_{i,l},&
\rho_{i,l}&=\underline\rho_{i,l}+\Delta\rho_{i,l}\Theta_{i,l},\\
\Delta M_{t,i,l}&=\overline M_{t,i,l}-\underline M_{t,i,l},&
M_{t,i,l}&=\underline M_{t,i,l}+\Delta M_{t,i,l}\kappa_{t,i,l},\\
\Delta V_{i,j,l}&=\overline V_{i,j,l}-\underline V_{i,j,l},&
V_{i,j,l}&=\underline V_{i,j,l}+\Delta V_{i,j,l}\varrho_{i,j,l}.
\end{align*}
Here $t\in\mathcal T$, $i\in I$, $l\in L$, and $(i,j)\in\mathcal A$ as applicable.
The demand selector is defined as one at the upper endpoint and zero at the lower
endpoint; the other selectors follow the same construction. No separate rule for
coincident endpoints is supplied.

The following are the authors' meanings of the product auxiliaries, not additional
nonlinear equalities imposed alongside their replacement rows (C.3):
\begin{align*}
k_{t,i,l}&=q_{i,l}\beta_{t,i},&
a'_{t,i,l}&=a_{t,i,l}q_{i,l},&
h_{t,i,l}&=a'_{t,i,l}\beta_{t,i},\\
\Theta'_{i,l}&=\Theta_{i,l}q_{i,l},&
g_{t,i,l}&=\Theta'_{i,l}\beta_{t,i},&
F_{t,i,l}&=q_{i,l}\Gamma_{t,i},\\
\kappa'_{t,i,l}&=\kappa_{t,i,l}q_{i,l},&
\pi_{t,i,l}&=\kappa'_{t,i,l}\Gamma_{t,i},&
\eta_{i,j,l}&=q_{i,l}\psi_{i,j},\\
\varrho'_{i,j,l}&=\varrho_{i,j,l}q_{i,l},&
\Phi_{i,j,l}&=\varrho'_{i,j,l}\psi_{i,j},&
\varpi_{i,j,l}&=q_{j,l}\varphi_{i,j},\\
b_{i,j,l}&=\varrho_{i,j,l}q_{j,l},&
\Lambda_{i,j,l}&=b_{i,j,l}\varphi_{i,j}.&&
\end{align*}
The node/item families are indexed by $t\in\mathcal T$, $i\in I$, $l\in L$,
with $t$ absent for $\Theta'$; the arc families are indexed by
$(i,j)\in\mathcal A$, $l\in L$.

\subsection*{Complete linearization rows used in (9b)}
The lower and upper decorations on dual coefficients below denote fixed bounds,
not additional optimization variables. The source prints these families without
individual quantifiers; their component ranges are those just specified. The source prints inconsistent $\eta_{t,i,l}$ subscripts in (C.3i). The operative candidate selects the explicit arc-product definition on the same source page; the original conflicting subscripts remain documented below.
\begin{align}
&\begin{aligned}
k_{t,i,l}&\geq q_{i,l}\underline\beta_{t,i},&
k_{t,i,l}&\geq\beta_{t,i}+\overline\beta_{t,i}(q_{i,l}-1),\\
k_{t,i,l}&\leq q_{i,l}\overline\beta_{t,i},&
k_{t,i,l}&\leq\beta_{t,i}+\underline\beta_{t,i}(q_{i,l}-1)
\end{aligned}\tag{C.3a}\\
& a'_{t,i,l}\geq0,\quad a'_{t,i,l}\leq q_{i,l},\quad
 a'_{t,i,l}\leq a_{t,i,l},\quad a'_{t,i,l}\geq a_{t,i,l}+q_{i,l}-1
\tag{C.3b}\\
&\begin{aligned}
h_{t,i,l}&\geq a'_{t,i,l}\underline\beta_{t,i},&
h_{t,i,l}&\geq\beta_{t,i},{}+\overline\beta_{t,i}(a'_{t,i,l}-1),\\
h_{t,i,l}&\leq a'_{t,i,l}\overline\beta_{t,i},&
h_{t,i,l}&\leq\beta_{t,i}+\underline\beta_{t,i}(a'_{t,i,l}-1)
\end{aligned}\tag{C.3c}\\
&\Theta'_{i,l}\geq0,\quad\Theta'_{i,l}\leq q_{i,l},\quad
\Theta'_{i,l}\leq\Theta_{i,l},\quad\Theta'_{i,l}\geq\Theta_{i,l}+q_{i,l}-1
\tag{C.3d}\\
&\begin{aligned}
g_{t,i,l}&\geq\Theta'_{i,l}\underline\beta_{t,i},&
g_{t,i,l}&\geq\beta_{t,i},{}+\overline\beta_{t,i}(\Theta'_{i,l}-1),\\
g_{t,i,l}&\leq\Theta'_{i,l}\overline\beta_{t,i},&
g_{t,i,l}&\leq\beta_{t,i}+\underline\beta_{t,i}(\Theta'_{i,l}-1)
\end{aligned}\tag{C.3e}\\
&\begin{aligned}
F_{t,i,l}&\geq q_{i,l}\underline\Gamma_{t,i},&F_{t,i,l}&\geq\Gamma_{t,i},\\
F_{t,i,l}&\leq0,&F_{t,i,l}&\leq\Gamma_{t,i}+\underline\Gamma_{t,i}(q_{i,l}-1)
\end{aligned}\tag{C.3f}\\
&\kappa'_{t,i,l}\geq0,\quad\kappa'_{t,i,l}\leq q_{i,l},\quad
\kappa'_{t,i,l}\leq\kappa_{t,i,l},\quad
\kappa'_{t,i,l}\geq\kappa_{t,i,l}+q_{i,l}-1
\tag{C.3g}\\
&\begin{aligned}
\pi_{t,i,l}&\geq\kappa'_{t,i,l}\underline\Gamma_{t,i},&\pi_{t,i,l}&\geq\Gamma_{t,i},\\
\pi_{t,i,l}&\leq0,&\pi_{t,i,l}&\leq\Gamma_{t,i}+\underline\Gamma_{t,i}(\kappa'_{t,i,l}-1)
\end{aligned}\tag{C.3h}\\
&\begin{aligned}
\eta_{i,j,l}&\geq q_{i,l}\underline\psi_{i,j},&\eta_{i,j,l}&\geq\psi_{i,j},\\
\eta_{i,j,l}&\leq0,&\eta_{i,j,l}&\leq\psi_{i,j}+\underline\psi_{i,j}(q_{i,l}-1)
\end{aligned}\tag{C.3i}\\
&\varrho'_{i,j,l}\geq0,\quad\varrho'_{i,j,l}\leq q_{i,l},\quad
\varrho'_{i,j,l}\leq\varrho_{i,j,l},\quad
\varrho'_{i,j,l}\geq\varrho_{i,j,l}+q_{i,l}-1
\tag{C.3j}\\
&\begin{aligned}
\Phi_{i,j,l}&\geq\varrho'_{i,j,l}\underline\psi_{i,j},&\Phi_{i,j,l}&\geq\psi_{i,j},\\
\Phi_{i,j,l}&\leq0,&\Phi_{i,j,l}&\leq\psi_{i,j}+\underline\psi_{i,j}(\varrho'_{i,j,l}-1)
\end{aligned}\tag{C.3k}\\
&\begin{aligned}
\varpi_{i,j,l}&\geq q_{j,l}\underline\varphi_{i,j},&\varpi_{i,j,l}&\geq\varphi_{i,j},\\
\varpi_{i,j,l}&\leq0,&\varpi_{i,j,l}&\leq\varphi_{i,j}+\underline\varphi_{i,j}(q_{j,l}-1)
\end{aligned}\tag{C.3l}\\
&b_{i,j,l}\geq0,\quad b_{i,j,l}\leq q_{j,l},\quad
 b_{i,j,l}\leq\varrho_{i,j,l},\quad b_{i,j,l}\geq\varrho_{i,j,l}+q_{j,l}-1
\tag{C.3m}\\
&\begin{aligned}
\Lambda_{i,j,l}&\geq b_{i,j,l}\underline\varphi_{i,j},&\Lambda_{i,j,l}&\geq\varphi_{i,j},\\
\Lambda_{i,j,l}&\leq0,&\Lambda_{i,j,l}&\leq\varphi_{i,j}+\underline\varphi_{i,j}(b_{i,j,l}-1)
\end{aligned}\tag{C.3n}
\end{align}
\noindent\emph{Printed punctuation and index inconsistency.} The page image has a comma
between $\beta_{t,i}$ and the following plus sign in the second inequality of both
(C.3c) and (C.3e). In (C.3i), the last two inequalities use $\eta_{t,i,l}$ rather than
the arc-indexed $\eta_{i,j,l}$ defined in the preceding prose. This executable candidate selects that explicit source product form and binds all four C.3i rows to the same arc/level auxiliary. It does not claim equivalence to a literal reading that treats both subscript families as distinct. No additional nonlinear product equality is imposed alongside these four replacement rows.

\subsection*{Bound coefficients and the Appendix D discussion}
Appendix C identifies $\underline\beta,\overline\beta,
\underline\Gamma,\underline\psi,\underline\varphi$ as bound coefficients in the
McCormick inequalities. They are distinct from the unmarked dual variables.
Appendix D gives the following bound discussion (page 48):
\[
-c_t^h\leq\beta_{t,i}\leq c_t^u,\qquad
\underline\beta_{t,i}=-c_t^h,\qquad \overline\beta=c_t^u,
\qquad k_{t,i,l}\in[\underline\beta_{t,i},\overline\beta_{t,i}].
\]
The upper-bound definition is printed without component subscripts. For the procurement
dual, the discussion states
\[
\Gamma_{t,i}\leq c_t^p-\beta_{t,i},\qquad \Gamma_{t,i}\leq0,
\qquad \Gamma_{t,i}=\min\{c_t^p-\beta_{t,i},0\}
\quad\text{in an optimal solution},
\]
\[
\overline\Gamma_{t,i}=0,\qquad
\underline\Gamma_{t,i}=
\begin{cases}
c_t^p-c_t^u,&c_t^p\leq c_t^u,\\
0,&\text{otherwise}.
\end{cases}
\]
Its final paragraph starts from
\[
v_t\psi_{i,j}+v_t\varphi_{i,j}\leq c^t_{i,j}+\beta_{t,i}-\beta_{t,j},
\qquad \psi_{i,j}\leq0,\quad\varphi_{i,i}\leq0,
\]
and the supplied text gives the bound statement
\[
\underline\psi_{i,j}\leq
\min_t\left\{(1/v_t)(c^t_{i,j}-c_t^h-c_t^u\right\}
\quad\text{if }c^t_{i,j}\leq c_t^h+c_t^u,
\qquad \underline\psi_{i,j}=0\quad\text{otherwise}.
\]
This is retained as a bound discussion, not strengthened to an equality or repaired
into a different inequality. Its condition contains $t$ outside the minimum and
its parenthesis is incomplete in the supplied text. The page-48 image is not supplied,
so the decorations in this last bound statement, unlike the coefficients in the
page-47 inequalities, remain image-unverified. No formula assigning
$\underline\varphi_{i,j}$ appears in the supplied Appendix D text; it remains the
abstract bound coefficient used in (C.3l) and (C.3n). No numerical big-$M$, additional
positivity assumption on $v_t$, or symmetric repair for $\underline\varphi$ is added.

\subsection*{Recourse definition underlying the epigraph}
For feasible first-stage decisions and a given realization $\xi$, the recourse value
used in the ambiguity model is the following problem. These rows define $Q$ and the
origin of (7); they are not an additional primal block imposed on the variables of
(9).
\begin{equation}
Q(\boldsymbol o,\boldsymbol z,\xi):=
\min_{\boldsymbol y,\boldsymbol u,\boldsymbol e,\boldsymbol x}
\left\{\sum_{i\in I}\sum_{t\in\mathcal T}
(c_t^p y_{t,i}+c_t^u u_{t,i}+c_t^h e_{t,i})
+\sum_{t\in\mathcal T}\sum_{(i,j)\in\mathcal A}c^t_{i,j}x^t_{i,j}\right\}.
\tag{2a}
\end{equation}
\begin{align}
&\begin{aligned}
\text{s.t.}\quad
\sum_{j:(j,i)\in\mathcal A}x^t_{j,i}
-\sum_{j:(i,j)\in\mathcal A}x^t_{i,j}
+y_{t,i}-e_{t,i}+u_{t,i}
&=\sum_{l\in L}q_{i,l}(d_{t,i,l}-\rho_{i,l}z_{t,i})\\
&\quad-\left(1-\sum_{l\in L}q_{i,l}\right)z_{t,i},
\qquad\forall t\in\mathcal T,\ i\in I,
\end{aligned}\tag{2b}\\
&y_{t,i}\leq\sum_{l\in L}q_{i,l}M_{t,i,l},
\qquad\forall t\in\mathcal T,\ i\in I,\tag{2c}\\
&\sum_{t\in\mathcal T}v_t x^t_{i,j}
\leq\sum_{l\in L}q_{i,l}V_{i,j,l}
+\left(1-\sum_{l\in L}q_{i,l}\right)\hat V_{i,j},
\qquad\forall(i,j)\in\mathcal A,\tag{2d}\\
&\sum_{t\in\mathcal T}v_t x^t_{i,j}
\leq\sum_{l'\in L}q_{j,l'}V_{i,j,l'}
+\left(1-\sum_{l'\in L}q_{j,l'}\right)\hat V_{i,j},
\qquad\forall(i,j)\in\mathcal A,\tag{2e}\\
&y_{t,i}\geq0,\quad e_{t,i}\geq0,\quad u_{t,i}\geq0,\quad x^t_{i,j}\geq0,
\qquad\forall t\in\mathcal T,\ i\in(i,j)\in\mathcal A.\tag{2f}
\end{align}
The variables $y_{t,i},e_{t,i},u_{t,i}$ are continuous quantities procured, left unused,
and unmet, respectively; the defining prose gives them for all $t\in\mathcal T$,
$i\in\mathcal N$. The continuous variable $x^t_{i,j}$ is shipment of item $t$ through
arc $(i,j)\in\mathcal A$. All four are nonnegative by their prose definitions. The
malformed quantifier $i\in(i,j)\in\mathcal A$ in (2f) is retained.

If $\sum_{l\in L}q_{i,l}=0$, all inventory prepositioned at $i$ remains usable.
The accompanying arc-capacity interpretation is: nominal capacity $\hat V_{i,j}$
when neither endpoint is hit, $V_{i,j,l}$ or $V_{i,j,l'}$ when only $i$ or only $j$
is hit, respectively, and
\[
\min\{V_{i,j,l},V_{i,j,l'}\}
\]
when the two endpoints are hit at levels $l$ and $l'$. These are the authors'
interpretations of (2d)--(2e), not additional capacity inequalities.

The moment duality in Appendix B and the endpoint/product replacements in Appendix C
lead to the selected (9)--(10). The intermediate min--max formulations and algorithmic
master problems are not simultaneously enforced here. Proposition 3 states that, for
fixed outer arguments,
\[
H(\boldsymbol z,\boldsymbol\alpha,\boldsymbol\phi,
\boldsymbol\gamma,\boldsymbol\lambda,\boldsymbol\tau)<\infty,
\]
and that $H$ is convex piecewise linear in those arguments with finitely many pieces.

\paragraph{Source.}
K. S. Shehadeh and E. L. Tucker, \emph{Stochastic Optimization Models for Location and
Inventory Prepositioning of Disaster Relief Supplies}, arXiv:2012.05387v3,
physical PDF pages 10--17 and Appendices B--E, pages 45--49. Equations (3)--(4) and the
Appendix B and D passages are available in the supplied full-page text; their original
page images are not supplied. The supplied images of pages 12, 16, 17, and 47 establish
the displayed typographical issues and the bound accents in the linearization rows.
\end{document}